\section{The \repay{} algorithm}\label{sec:algorithm}

\repay{} (\textsc{re}al-time \textsc{p}riority with el\textsc{a}stic pa\textsc{y}back) is a drift-plus-penalty rule with two families of virtual queues.
Three \emph{tier queues} track the business constraints and persist as devices come and go.
In addition, every elastic device carries, for the time it is present, one \emph{deficit queue} per RT class that records its outstanding class envy.

\subsection{State}

\paragraph{Tier queues.}
For $j\in\{PS,YS,FF\}$, $Q_j(1)=0$ and
\begin{equation}\label{eq:Q}
Q_j(t+1)=\max\big(Q_j(t)+e_{j,t},\,0\big).
\end{equation}
They induce a per-unit \emph{tier price} $\beta_{P,t}=Q_{PS}(t)$, $\beta_{Y,t}=Q_{YS}(t)$, $\beta_{F,t}=Q_{FF}(t)$: the bonus that a unit of rate delivered to the tier earns in slot $t$.

\paragraph{Deficit queues.}
When an elastic device $i$ arrives, $D_{i,k}=0$ for every $k\in\cK$. While it is present,
\begin{equation}\label{eq:D}
D_{i,k}(t+1)=\max\big(D_{i,k}(t)+r_{i,k,t}-x_{i,t},\,0\big),
\end{equation}
with the reference $r_{i,k,t}$ of \eqref{eq:ref}. The deficits are discarded when the device leaves.
Unrolling the recursion gives $D_{i,k}(t+1)=\max_{s\le t}\sum_{\tau=s}^{t}(r_{i,k,\tau}-x_{i,\tau})$ with the maximum taken over windows that start inside the sojourn (and the empty window): the deficit is exactly the largest class envy over all windows that end now.

\subsection{Per-slot problem}

The reference $a\min(z,d_i)$ is not linear in the decision, and putting it directly into the objective would destroy concavity.
We replace $z\mapsto\min(z,d)$ on $[\rmin_k,\rmax_k]$ by the affine majorant
\begin{equation}\label{eq:maj}
m_d(z)=\sigma+s\,z,\qquad
(s,\sigma)=\begin{cases}
(1,0), & d\ge\rmax_k,\\
(0,d), & d\le\rmin_k,\\
\Big(\dfrac{d-\rmin_k}{\rmax_k-\rmin_k},\ (1-s)\,d\Big), & \text{otherwise,}
\end{cases}
\end{equation}
which satisfies $m_d(z)\ge\min(z,d)$ on the interval, with equality whenever $d\notin(\rmin_k,\rmax_k)$.
For device $i$ and class $k$ we write $(s_{i,k},\sigma_{i,k})$ for the coefficients obtained with $d=d_{i,t}$.
Group $(g,k)$ then faces the \emph{debt coefficients}
\begin{equation}\label{eq:debt}
\pi_{g,k}=\sum_{i\in\cE_t:\,g_i=g}D_{i,k}\,s_{i,k},\qquad
\zeta_{g,k}=\sum_{i\in\cE_t:\,g_i=g}D_{i,k}\,\sigma_{i,k},
\end{equation}
and device $i$ earns the \emph{bonus} $D_i=\sum_{k\in\cK}D_{i,k}$ per unit of rate.

In slot $t$, \repay{} solves
\begin{equation}\label{eq:slot}
\begin{aligned}
\max_{a,y,x}\quad &\sum_{(g,k):\,n_{g,k}\ge 1}\Big[V w_g\,n_{g,k}\,a_{g,k}\,u_k\!\big(y_{g,k}/a_{g,k}\big)+\big(n_{g,k}\beta_g-\pi_{g,k}\big)y_{g,k}-\zeta_{g,k}\,a_{g,k}\Big]\\
&+\sum_{i\in\cE_t}\Big[V w_{g_i}\,u(x_i)+\big(\beta_{g_i}+D_i\big)x_i\Big]\\
\text{s.t.}\quad &\textstyle\sum_{g,k}n_{g,k}y_{g,k}+\sum_i x_i\le C_t,\qquad 0\le x_i\le d_{i,t},\\
&0\le a_{g,k}\le 1,\qquad a_{g,k}\rmin_k\le y_{g,k}\le a_{g,k}\rmax_k ,
\end{aligned}
\end{equation}
where $a\,u_k(y/a)$ is taken to be $0$ at $a=0$ and $V>0$ is the usual drift-plus-penalty parameter.
Up to terms that do not depend on the decision, the objective equals
\begin{equation}\label{eq:dpp}
V\,U_t-\sum_jQ_j(t)\,e_{j,t}-\sum_{i\in\cE_t}\sum_{k\in\cK}D_{i,k}(t)\,\big(a_{g_i,k}\,m_{d_{i,t}}(z_{g_i,k})-x_i\big).
\end{equation}

The three groups of terms have a direct reading.
An RT group pays, per unit of rate, the sum of the deficits it has built up among its elastic tier-mates: priority is bought on credit, and the price of further priority rises with the outstanding debt.
An elastic device that is owed volume bids $D_i$ above its marginal utility, so it is served before everything else once capacity returns.
Tier prices shift capacity between tiers but, being per unit of rate, do not change the comparison between classes inside a tier.

\begin{lemma}\label{lem:concave}
Problem \eqref{eq:slot} is a concave maximisation over a polyhedron.
\end{lemma}
\begin{proof}
$(a,y)\mapsto a\,u_k(y/a)$ is the perspective of the concave function $u_k$, hence jointly concave on $a>0$, and its closure takes the value $0$ at the origin.
The debt terms are linear in $(a,y)$ because $a\,m_d(y/a)=\sigma a+s\,y$.
The elastic terms are concave, and all constraints are linear.
\end{proof}

Writing the RT decision as $(a,z)$ would not work: $a\,u_k(z)$ is not jointly concave, and neither is $-D\min(az,d)$.
The perspective variables and the majorant are what make the per-slot problem tractable.

\subsection{Exact solution by bisection}

Let $\lambda\ge 0$ be the multiplier of the capacity constraint.
For a given $\lambda$ the problem separates.

\emph{Elastic device.} With $p_i=\lambda-\beta_{g_i}-D_i$,
\begin{equation}\label{eq:xresp}
x_i(\lambda)=\begin{cases}d_i, & p_i\le 0,\\
\min\big\{d_i,\ \max\big(0,\ V w_{g_i}/p_i-1\big)\big\}, & p_i>0 .\end{cases}
\end{equation}

\emph{RT group.} For fixed $z$ the group's Lagrangian is linear in $a$, equal to $a\,\psi_{g,k}(z;\lambda)$ with
\begin{equation}\label{eq:psi}
\psi_{g,k}(z;\lambda)=n_{g,k}\big[V w_g\,u_k(z)+(\beta_g-\lambda)z\big]-\pi_{g,k}z-\zeta_{g,k}.
\end{equation}
Hence $z^*(\lambda)=\arg\max_{z\in[\rmin_k,\rmax_k]}\psi_{g,k}(z;\lambda)$ and $a^*(\lambda)=\ind\{\psi_{g,k}(z^*;\lambda)>0\}$.
For video, with $q=n_{g,k}(\lambda-\beta_g)+\pi_{g,k}$, $z^*=\rmax_V$ if $q\le 0$ and $z^*=\min\{\rmax_V,\max(\rmin_V,\,n_{g,k}Vw_g/q)\}$ otherwise; for gaming $z^*=r_G$.

The load $S(\lambda)=\sum n\,a^*z^*+\sum_i x_i(\lambda)$ is nonincreasing in $\lambda$.
If $S(0)\le C_t$ the solution is the one at $\lambda=0$.
Otherwise we bisect on $\lambda$ until $S(\lambda_{lo})>C_t\ge S(\lambda_{hi})$ with $\lambda_{hi}-\lambda_{lo}$ below a tolerance, and take the allocation at $\lambda_{hi}$.
A group that is admitted at $\lambda_{lo}$ and blocked at $\lambda_{hi}$ is \emph{marginal}: its $\psi$ vanishes at the optimal multiplier, so any $a\in[0,1]$ is optimal for it, and we give it the fraction that fills the remaining capacity.

\begin{proposition}[Solver]\label{prop:solver}
As the tolerance tends to zero, the procedure returns an optimal solution of \eqref{eq:slot}.
One slot costs $O\big((N_t+|\cK|\cdot 3)\log(1/\epsilon)\big)$ operations for $N_t$ active devices and relative tolerance $\epsilon$.
\end{proposition}
\begin{proof}
By Lemma~\ref{lem:concave} and Slater's condition (the zero allocation leaves capacity $C_t>0$ unused) strong duality holds with some $\lambda^*\ge0$, and a feasible point is optimal if and only if it maximises the Lagrangian at $\lambda^*$ and satisfies complementary slackness.
The Lagrangian maximisers $x_i(\lambda)$ and $z^*(\lambda)$ are unique and continuous in $\lambda$ because $u$ and $u_V$ are strictly concave (and $z$ is fixed for gaming).
The admission variable is free in $[0,1]$ exactly for groups with $\psi(z^*;\lambda^*)=0$.
$S$ is nonincreasing: $x_i$ decreases in $\lambda$; $\psi$ has decreasing differences in $(z,\lambda)$, so $z^*$ decreases; and $\max_z\psi$ decreases, so $a^*$ decreases.
Therefore either $S(0)\le C_t$ and $\lambda^*=0$, or $S(\lambda^*-)\ge C_t\ge S(\lambda^*+)$ and capacity can be met exactly by choosing the free admission fractions; this is what the procedure constructs as $\lambda_{hi}\downarrow\lambda^*$.
Each evaluation of $S$ is linear in the number of devices and groups.
\end{proof}

\subsection{From fractions to flows}

Only marginal groups have a fractional $a$, generically at most one per slot.
For such a group \repay{} admits $M=\lfloor n a\rfloor$ flows at rate $z$, giving precedence to flows that were served in the previous slot so that ongoing calls are not interrupted, and then to earlier arrivals.
The leftover capacity $(na-M)z$ serves one more flow if it is at least $\rmin_k$; otherwise it is handed to elastic devices that are below their cap, in decreasing order of marginal value.
Rounding therefore never reduces an elastic allocation.
The deficit update \eqref{eq:D} uses the fluid reference $a\min(z,d_i)$, which is at least the average rate of the admitted flows of the group and so errs in favour of the elastic device.

\begin{algorithm}[t]
\caption{\repay{}, slot $t$}\label{alg:repay}
\begin{algorithmic}[1]
\State observe $C_t$, the groups $n_{g,k}$, and for each elastic device its tier and verified demand cap $d_i$
\State $\beta_g\gets$ tier prices from $Q$; \ $(s_{i,k},\sigma_{i,k})\gets$ majorant coefficients \eqref{eq:maj}; \ $\pi,\zeta\gets$ \eqref{eq:debt}; \ $D_i\gets\sum_kD_{i,k}$
\State $(a,z,x,\lambda)\gets$ solve \eqref{eq:slot} by bisection on $\lambda$
\State round admissions to whole flows (incumbents first) and serve
\ForAll{elastic $i$, class $k$} $D_{i,k}\gets\max\big(D_{i,k}+a_{g_i,k}\min(z_{g_i,k},d_i)-x_i,\ 0\big)$ \EndFor
\ForAll{$j$} $Q_j\gets\max(Q_j+e_j,\ 0)$ \EndFor
\State arrivals start with $D=0$; departures drop their deficits; a device that changes its declared class keeps the new class for at least $\tau$ slots
\end{algorithmic}
\end{algorithm}

The rule has four knobs: $V$, which trades utility against the size of the envy bound; the tier weights $w_g$; the commercial targets $(\theta_P,\theta_Y,\phi)$; and the \emph{hold-down} time $\tau$ that a subscriber must keep a declared class before changing it.
Section~\ref{sec:analysis} explains why the hold-down is needed.

\subsection{Variants}\label{sec:variants}

We also evaluate three modifications.
\textsc{Repay-E} (entrant offer) defines a reference even when the group $(g,k)$ is empty: the offer that a single new class-$k$ flow of tier $g$ would receive at the realised $\lambda$, obtained from \eqref{eq:psi} with $n=1$. It closes a loophole in thin markets, where a lone elastic device could become the only RT flow of its tier.
\textsc{Repay-WC} (work-conserving) hands capacity that is left unused at $\lambda=0$ to RT groups, which the debt price may otherwise keep below $\rmax_k$ for a few slots after a congestion episode.
\textsc{No-debt} drops the debt coefficients ($\pi=\zeta=0$) and keeps only the elastic bonus.
The deterministic guarantees of Section~\ref{sec:envy} hold for all three; the optimality guarantee of Section~\ref{sec:opt} is proved for \repay{} itself.
